%\ifodd\value{page}\hbox{}\newpage\fi
\chapter*{Śrī Lalitā Sahasranāma — Śloka 13}
\addcontentsline{toc}{chapter}{Śloka 13 - Nomi 28,29}

\begin{sloka}
{\sanskritfont
\begin{tabular}{c}
कनकाङ्गदकेयूरकमनीयभुजान्विता ।\\
रत्नग्रैवेयचिन्ताकलोलमुक्ताफलान्विता ॥ १३॥
\end{tabular}

\vspace{1.5em}

{\middlesize \iastfont
kanakāṅgada-keyūra-kamanīya-bhujānvitā |\\
ratna-graiveya-cintāka-lola-muktā-phalānvitā ||13||}

\vspace{1em}

{\middlesize
\anvaya{%
कनकाङ्गदकेयूरकमनीयभुजान्विता ।
रत्नग्रैवेयचिन्ताकलोलमुक्ताफलान्विता ।
}
}

\middlesize \iastfont \itmean{%
«Adornata di bracciali d’oro e di ornamenti alle braccia, dalle braccia di grazia incantevole;  
adornata di collane di gemme e di perle mobili che oscillano sul petto.»%
}

}
\end{sloka}

\wordmeaning{%
\wordline{कनक-अङ्गद-केयूर-कमनीय-भुज-अन्विता}
{kanaka-aṅgada-keyūra-kamanīya-bhuja-anvitā}
{ adornata (\textit{anvitā}) di bracciali d’oro (\textit{kanaka-aṅgada}) e di ornamenti per le braccia (\textit{keyūra}), dalle braccia (\textit{bhuja}) graziose e incantevoli (\textit{kamanīya}); }%

\wordline{रत्न-ग्रैवेय-चिन्ताक-लोल-मुक्ता-फल-अन्विता}
{ratna-graiveya-cintāka-lola-\\muktā-phala-anvitā}
{ adornata (\textit{anvitā}) di collane (\textit{graiveya}) di gemme (\textit{ratna}), con pendenti (\textit{cintāka}) e perle (\textit{muktā-phala}) mobili e oscillanti (\textit{lola}); }%
}

\textit{Śloka} 13 sposta lo sguardo sulle braccia e sul petto di \textit{Lalitā}, introducendo un linguaggio di ornamento che non è mera decorazione, ma manifestazione visibile di grazia e pienezza. I bracciali e gli ornamenti alle braccia non aggiungono bellezza: la seguono, la accompagnano, come se la forma stessa li richiedesse.

Nel composto \textit{bhujānvitā}, il termine \textit{bhuja} designa l’intero arto; alcune letture tradizionali vi riconoscono anche un’allusione alle quattro braccia della Dea come segno di protezione e potenza. Qui l’accento resta sulla grazia e sull’ornamento, secondo una prospettiva prevalentemente contemplativa.

Nel secondo \textit{pāda}, le collane di gemme e le perle oscillanti introducono il tema del movimento. Le perle non sono immobili: vibrano, riflettono la luce, partecipano al ritmo del corpo. La bellezza qui non è statica, ma viva, ritmica, continuamente rinnovata.

Alcuni commentari leggono \textit{muktā-phala} non solo come “perle”, ma anche come “frutti” (\textit{phala}) donati dalla Dea ai devoti, suggerendo una sovrapposizione simbolica tra ornamento visibile e beneficio spirituale. In questa sede la resa resta sul piano letterale, senza escludere tale risonanza ulteriore.

Questo \textit{śloka} contiene due \textit{nāma}.  
\textbf{\textit{Kanakāṅgada-keyūra-kamanīya-bhujānvitā}} descrive le braccia ornate come luogo di grazia naturale.  
\textbf{\textit{Ratna-graiveya-cintāka-lola-muktā-phalānvitā}} presenta gli ornamenti del petto come segni di una bellezza in movimento, luminosa e vitale.
